\section{Tarefa 5.4: Projeto e montagem de um somador completo de 1 bit}

Esta tarefa consistiu no projeto de um somador completo de 1 bit com propagacao do vai-um. A celula basica deve produzir a soma $S_i$ e o transporte $C_{i-1}$ a partir das entradas $A_i$, $B_i$ e $C_i$, com implementacao final usando somente portas NAND de duas entradas.

\subsection{Tabela-verdade}

\begin{table}[H]
    \centering
    \caption{Tabela-verdade da celula do somador completo.}
    \begin{tabular}{ccccc}
        \toprule
        $A_i$ & $B_i$ & $C_i$ & $S_i$ & $C_{i-1}$ \\
        \midrule
        0 & 0 & 0 & & \\
        0 & 0 & 1 & & \\
        0 & 1 & 0 & & \\
        0 & 1 & 1 & & \\
        1 & 0 & 0 & & \\
        1 & 0 & 1 & & \\
        1 & 1 & 0 & & \\
        1 & 1 & 1 & & \\
        \bottomrule
    \end{tabular}
\end{table}

\subsection{Expressoes logicas}

\[
S_i = \underline{\hspace{8cm}}
\]

\[
C_{i-1} = \underline{\hspace{8cm}}
\]

\subsection{Implementacao usando NAND2}

Apresente aqui a minimizacao das expressoes e a implementacao usando somente portas NAND de duas entradas.

\begin{figure}[H]
    \centering
    % \includegraphics[width=0.8\textwidth]{questao4/somador_nand.png}
    \fbox{\rule{0pt}{5cm}\rule{0.8\textwidth}{0pt}}
    \caption{Espaco reservado para o circuito em NAND2 do somador completo.}
\end{figure}

\subsection{Resultados experimentais}

\begin{table}[H]
    \centering
    \caption{Resultados experimentais do somador completo de 1 bit.}
    \begin{tabular}{ccccc}
        \toprule
        $A_i$ & $B_i$ & $C_i$ & $S_i$ & $C_{i-1}$ \\
        \midrule
        0 & 0 & 0 & & \\
        0 & 0 & 1 & & \\
        0 & 1 & 0 & & \\
        0 & 1 & 1 & & \\
        1 & 0 & 0 & & \\
        1 & 0 & 1 & & \\
        1 & 1 & 0 & & \\
        1 & 1 & 1 & & \\
        \bottomrule
    \end{tabular}
\end{table}

Discuta aqui a coerencia entre os resultados observados e o comportamento esperado.
